% Figure: temperature-entropy diagram with the vapour dome and a
% schematic Rankine-style cycle inscribed. The dome peak is the critical
% point; isobars run flat across the dome. Cycle: pump (left limb up),
% boiler (flat across to superheat), turbine (down to two-phase),
% condenser (flat back). Schematic coordinates.
\begin{figure}[htbp]
\centering
\definecolor{engteal}{RGB}{30,120,120}
\begin{tikzpicture}
\begin{axis}[
  width=12cm, height=8cm,
  xlabel={specific entropy $s$ (arb.\ units)},
  ylabel={temperature $T$ (arb.\ units)},
  xmin=0, xmax=10, ymin=0, ymax=8,
  axis lines=left,
  xtick=\empty, ytick=\empty,
  tick label style={font=\scriptsize},
  label style={font=\small},
  clip=true,
]
% vapour dome
\addplot[engblue, very thick, smooth] coordinates {
  (1.6,1.0) (1.9,2.4) (2.4,3.8) (3.2,5.0) (4.2,5.8) (5.0,6.2)
};
\addplot[engblue, very thick, smooth] coordinates {
  (5.0,6.2) (6.0,5.9) (6.9,5.2) (7.7,4.2) (8.4,2.8) (8.9,1.0)
};
\addplot[only marks, mark=*, engblack, mark size=2.0pt] coordinates {(5.0,6.2)};
\node[engblack, font=\scriptsize, anchor=south] at (axis cs:5.0,6.4) {critical point};
\node[engblue, font=\scriptsize] at (axis cs:3.2,2.4) {two-phase};
% Rankine-style cycle:
% 1: saturated liquid at low T (2.0,1.4) ; 2: compressed liquid after pump (2.05,1.7)
% boiler heats liquid up left limb to saturation, then across to superheat
% 3: superheated steam (7.0,5.2) ; turbine expands to 4 (6.6,1.6) two-phase
% condenser back to 1.
\addplot[engred, very thick] coordinates {(2.0,1.4) (2.05,1.8)}; % pump
% boiler: up the saturated-liquid limb then to superheat point
\addplot[engorange, very thick, smooth] coordinates {
  (2.05,1.8) (2.4,3.0) (3.0,4.2) (4.0,5.4) (5.0,6.2) (6.2,6.6) (7.0,6.8)
};
% turbine: expand down into the two-phase region
\addplot[engred, very thick] coordinates {(7.0,6.8) (7.0,1.4)};
% condenser: flat back to state 1
\addplot[engteal, very thick] coordinates {(7.0,1.4) (2.0,1.4)};
% state labels
\addplot[only marks, mark=*, engblack, mark size=1.4pt] coordinates {(2.0,1.4) (7.0,6.8) (7.0,1.4)};
\node[engblack, font=\scriptsize, anchor=north east] at (axis cs:2.0,1.4) {1};
\node[engblack, font=\scriptsize, anchor=south west] at (axis cs:7.0,6.8) {3};
\node[engblack, font=\scriptsize, anchor=north west] at (axis cs:7.0,1.4) {4};
\node[engorange, font=\scriptsize, anchor=south] at (axis cs:4.6,6.7) {boiler};
\node[engred, font=\scriptsize, anchor=west] at (axis cs:7.1,4.1) {turbine};
\node[engteal, font=\scriptsize, anchor=north] at (axis cs:4.5,1.4) {condenser};
\end{axis}
\end{tikzpicture}
\caption[Temperature-entropy diagram with an inscribed cycle]{The
temperature-entropy diagram with the vapour dome (blue) and a schematic
steam-power cycle inscribed. Saturated liquid lies along the left limb of
the dome, saturated vapour along the right limb, and the peak is the
critical point. The cycle runs: a pump raises the pressure of the
saturated liquid (state~1), the boiler adds heat along the
constant-pressure path through the two-phase region and into the
superheated region (state~3), the turbine expands the steam back into the
two-phase region (state~4), and the condenser rejects heat at constant
pressure back to state~1. The area enclosed by the cycle on the $T$-$s$
diagram equals the net heat, and the clockwise direction marks it as a
work-producing engine. The same diagram traversed counter-clockwise is a
refrigeration cycle. Coordinates are schematic.}
\label{fig:vol04ch02:ts-dome-cycle}
\end{figure}
